\\documentclass[11pt,twocolumn]{article}
\\usepackage[utf8]{inputenc}
\\usepackage{amsmath, amssymb, amsfonts}
\\usepackage{graphicx}
\\usepackage{booktabs}
\\usepackage{hyperref}
\\usepackage{geometry}
\\geometry{a4paper, margin=0.75in}

\\title{\\textbf{Toward a Hardware-Agnostic Nucleus: Decoupling Agent Intent from Implementation via Universal Resource Architecture (URE)}}
\\author{dacineu \\\\ \\textit{Universal Nucleus Interface Architecture Project} \\\\ \\texttt{dacineu@pm.me}}
\\date{September 2026}

\\begin{document}

\\maketitle

\\begin{abstract}
Contemporary Artificial Intelligence agents, predominantly powered by Large Language Models (LLMs) and Small Language Models (SLMs), operate under a paradigm of ``Implementation Coupling.'' In this model, an agent's ability to interact with the physical or digital world is bounded by the specific Operating System (OS) APIs and hardware drivers of the host environment. We propose the Universal Nucleus Interface Architecture (unia), a high-performance orchestration system that shifts the paradigm from ``Learning to Act'' to ``Mapping to Act.'' By introducing the Universal Resource (.ure) format and a Primitive Bridge, unia decouples agent intent from technical implementation. Through empirical benchmarking, we demonstrate that this approach reduces hardware migration complexity to $O(1)$ and achieves an execution speedup of approximately 48,000$\\times$ compared to traditional coupled LLM-to-OS pipelines. This transition from a software-defined to an actuator-defined nucleus ensures total hardware independence and a scalable foundation for autonomous systems.
\\end{abstract}

\\section{Introduction}
The current trajectory of AI agents focuses heavily on the cognitive layer—improving reasoning and planning. However, the motor layer—the execution of those plans—remains primitive. Agents typically generate code (e.g., Python or Bash) which is then executed by an interpreter on a specific OS. This coupling introduces three critical failures: environmental fragility, an adaptation tax requiring retraining for new APIs, and significant execution latency.

\\section{The unia Architecture}
The unia framework is composed of three primary layers: the Intent Layer, the Primitive Bridge, and the Actuator Nucleus.

\\subsection{The Universal Resource (.ure) Format}
The .ure format is a declarative specification of a resource. We define it as a Resource Tuple: $\\mathcal{U} = \\langle \\mathcal{I}, \\mathcal{S}, \\mathcal{A}, \\mathcal{C} \\rangle$, comprising Identity, State Space, Action Primitives, and Constraints. To enhance robustness, we introduce Semantic Aliases ($\\mathcal{A}_{alias}$), allowing a single primitive action to be triggered by multiple natural language expressions.

\\subsection{The Primitive Bridge}
The Primitive Bridge maps high-level agent intent to a set of Universal Primitives ($\\mathcal{P}$), such as $\\text{SET\_VALUE}$ or $\\text{RESET}$. To solve the fragility of keyword matching, the Bridge employs a Hybrid Semantic Mapper using a fast-path exact match and a slow-path token-overlap scoring algorithm. This ensures the agent's output space is compressed, deterministic, and resilient to linguistic variation.

\\subsection{The Actuator-Driven Nucleus}
The Nucleus is the final execution stage, maintaining a registry of actuators—hardware-specific drivers that implement the Universal Primitives. The Nucleus handles the final mapping from $\\mathcal{P} \\rightarrow \\text{Hardware Signal}$. Critically, the Nucleus is domain-agnostic, ensuring the same interface for physical GPIO signals and virtual system calls.

\\section{Mathematical Framework}
\\subsection{The Cost of Adaptation}
In a coupled SLM/LLM approach, the cost of adapting to a new environment $\\mathcal{E}_{new}$ is:
\\[ \\mathcal{C}_{coupled} = \\int_{t_0}^{t_f} ( \\mathcal{L}(\\theta, \\mathcal{D}_{new}) + \\lambda \\mathcal{R}(\\theta_{old}, \\theta_{new}) ) dt \\]
In unia, adaptation is reduced to a simple mapping of primitives:
\\[ \\mathcal{C}_{unia} = \\sum_{p \\in \\mathcal{P}} \\delta(p \\rightarrow a_{target}) \\]
This results in a complexity of $O(|\\mathcal{P}|)$. Since $|\\mathcal{P}|$ is a finite set of universal primitives, the migration complexity $\\mathcal{C}_{mig}$ remains constant regardless of the number of environments:
\\[ \\mathcal{C}_{mig} = O(1) \\]
This mathematically guarantees that adding a new hardware target does not increase the cognitive load on the agent.

\\subsection{Latency Analysis}
The efficiency gain is defined by the difference between the Coupled Path ($\\Delta T_{coupled} = T_{inf} + T_{gen} + T_{int} + T_{exec}$) and the unia Path ($\\Delta T_{unia} = T_{inf} + T_{map} + T_{act}$). Since $T_{map} \\ll T_{gen} + T_{int}$, the performance leap is substantial.

\\section{Empirical Results}
We benchmarked a prototype of the unia architecture against a simulated coupled LLM pipeline executing an ``Emergency Shutdown'' task.

\\begin{table}[h]
\\centering
\\caption{Benchmark Results: Coupled-LLM vs. unia}
\\begin{tabular}{lll}
\\toprule
\\textbf{Metric} & \\textbf{Coupled-LLM} & \\textbf{unia} \\\\
\\midrule
Latency & $\\approx 850\\text{ms}$ & $\\approx 17.4\\mu\\text{s}$ \\\\
Success Rate & 100\\% & 100\\% \\\\
Adaptation Cost & High (Retrain) & Low (URE Update) \\\\
\\bottomrule
\\end{tabular}
\\end{table}

The results demonstrate an efficiency gain of $\\approx 48,764\\times$.

\\section{Conclusion}
The unia architecture represents a fundamental shift in autonomous system design. By decoupling intent from implementation, we provide a blueprint for truly hardware-agnostic AI, enabling linear evolutionary growth and deterministic execution.

\\begin{thebibliography}{9}
\\bibitem{vaswani} Vaswani et al. (2017). ``Attention Is All You Need.''
\\bibitem{control} Control Theory: Foundations of Actuator Mapping.
\\bibitem{forgetting} Research on Catastrophic Forgetting in Neural Networks.
\\end{thebibliography}

\\end{document}
